% Autorska priprava k Realisticky.cz 020318–020319; ulohy jsou autorske a obmenene.
% Zamerne nenahrazuje vyklad presnymi kopiemi cizich textu/prikladu.
\documentclass[11pt,a4paper]{article}
\input{../../sablony/pripravy/priprava.sty}
\usepackage{tikz}
\nastavPripravu{1. ročník}{Rovnice a nerovnice}{92}{Slovní úlohy vedoucí na soustavy}
\begin{document}
\blok{Návaznost a motivace}
Na soustavy navážeme úlohami z výroby, dopravy a směsí. Zásadní je propojit každou podmínku s jednou rovnicí a správně vyložit souřadnice dvojice řešení.
\blok{Princip a odvození}
Zaveď dvě neznámé, sestav dvě nezávislé rovnice, zvol metodu a vrať se k jednotkám. U směsí nezaměň hmotnost roztoku s hmotností složky.
\blok{Řešené příklady: od jednoduchého ke složitějšímu}
\begin{prikladReseniBox}[title={Příklad 1}]
Dvě dílny pracovaly dohromady 15 hodin a vyrobily 82 kusů. První vyrábí 6, druhá 4 kusy za hodinu. Jak dlouho pracovala každá?
\tcblower
\begin{align*}
x+y&=15 \\
 6x+4y&=82 \\
 6x+4(15-x)&=82 \\
 2x&=22 \\
 x&=11,\quad y=4
\end{align*}

Kontrola $11\cdot6+4\cdot4=82$.
\end{prikladReseniBox}
\begin{prikladReseniBox}[title={Příklad 2}]
Připravíme 10 kg směsi z 5% a 20% látky. Potřebujeme 1,1 kg účinné složky. Kolik kg které suroviny?
\tcblower
\begin{align*}
x+y&=10 \\
 0{,}05x+0{,}20y&=1{,}1 \\
 0{,}05(10-y)+0{,}20y&=1{,}1 \\
 0{,}15y&=0{,}6 \\
 y&=4,\quad x=6
\end{align*}

Skontroluj hmotnost účinné látky: $0{,}3+0{,}8=1{,}1$ kg.
\end{prikladReseniBox}
\poznamka{Zahrň kontrolu smyslu: u počtu autobusů nelze přijmout necelou hodnotu, i kdyby splňovala algebraický model.}
\blok{Samostatné procvičování}
\begin{zadaniReseniBox}[title={Úloha 1}]
50 krabic má dohromady 90 dílů; typ A obsahuje 3 a typ B 1 díl. Kolik kterých?
\tcblower
\begin{align*}
x+y&=50 \\
 3x+y&=90 \\
 2x&=40 \\
 x&=20,\quad y=30
\end{align*}
\end{zadaniReseniBox}
\begin{zadaniReseniBox}[title={Úloha 2}]
Pět autobusů má celkem 170 sedadel. Menší má 30 a větší 40. Kolik je kterých?
\tcblower
\begin{align*}
x+y&=5 \\
 30x+40y&=170 \\
 150+10y&=170 \\
 y&=2,\quad x=3
\end{align*}
\end{zadaniReseniBox}
\begin{zadaniReseniBox}[title={Úloha 3}]
Dvě suroviny stojí 60 a 100 Kč/kg. Potřebujeme 20 kg za 1600 Kč. Kolik které?
\tcblower
\begin{align*}
x+y&=20 \\
 60x+100y&=1600 \\
 1200+40y&=1600 \\
 y&=10,\quad x=10
\end{align*}
\end{zadaniReseniBox}
\blok{Volitelná rezerva}
\begin{zadaniReseniBox}[title={Úloha 1}]
Jaký počet 20Kč a 50Kč mincí dá dohromady 10 kusů a 350 Kč?
\tcblower
\begin{align*}
x+y&=10 \\
 20x+50y&=350 \\
 200+30y&=350 \\
 y&=5,\quad x=5
\end{align*}
\end{zadaniReseniBox}
\blok{Domácí úkol}
\begin{zadaniReseniBox}[title={Úloha 1}]
Dvě dílny vyrábí rychlostí 8 a 5 kg/h; pracují dohromady 12 hodin a vyrobí 81 kg. Doby?
\tcblower
\begin{align*}
x+y&=12 \\
 8x+5y&=81 \\
 8x+5(12-x)&=81 \\
 x&=7,\quad y=5
\end{align*}
\end{zadaniReseniBox}
\begin{zadaniReseniBox}[title={Úloha 2}]
Pět strojů typu A a B má dohromady výkon 26 kusů za hodinu. A má 6 a B 4 kusy za hodinu. Počty?
\tcblower
\begin{align*}
x+y&=5 \\
 6x+4y&=26 \\
 20+2x&=26 \\
 x&=3,\quad y=2
\end{align*}
\end{zadaniReseniBox}
\begin{zadaniReseniBox}[title={Úloha 3}]
Součet dvou věků je 53 let a rozdíl 7 let. Určete je.
\tcblower
\begin{align*}
x+y&=53 \\
 x-y&=7 \\
 2x&=60 \\
 x&=30,\quad y=23
\end{align*}
\end{zadaniReseniBox}
\poznamka{Orientační časové rozvržení: 5 + 12 + 12 + 10 + 6 = 45 min. Přednost má pochopení jádra; rezerva není součástí časového součtu. Domácí úlohy jsou určeny na samostatnou práci, řešení zůstávají učiteli.}
\end{document}
